% Part: proof-theory
% Chapter: normalization
% Section: normalization-thm

\documentclass[../../../include/open-logic-section]{subfiles}

\begin{document}

\olfileid[hi]{pt}{nor}{thm}

\olsection{प्रसामान्यीकरण प्रमेय}

!!^a{proof} \emph{प्रसामान्य} रूप में है यदि उसमें कोई कट खंड नहीं है। स्पष्ट
है कि ऐसा तभी और केवल तभी है जब उस पर न क्रम-विनिमय रूपांतरण लगाया जा सके, न
अपचयन रूपांतरण। किसी !!{proof} का \emph{प्रसामान्यीकरण} करने का अर्थ है कि
जब तक कोई कट शेष न रहे तब तक क्रम-विनिमय और अपचयन रूपांतरण लगाकर उसे
प्रसामान्य प्रमाण में बदल दिया जाए। यह सदा संभव है—यही \emph{प्रसामान्यीकरण
प्रमेय} है।

\begin{thm}\ollabel{thm:normalization} प्रत्येक \Log{N1i}-!!{proof}
$\delta$ का प्रसामान्यीकरण होता है; अर्थात् उसे कट-विहीन
!!a{proof}~$\delta'$ में बदला जा सकता है।
\end{thm}

\begin{proof}
  $\delta$ की कट कोटि और लंबाई पर आगमन से। आगमन परिकल्पना है कि कम कट कोटि,
  अथवा समान कट कोटि किंतु कम कट लंबाई वाले प्रत्येक !!{proof} का
  प्रसामान्यीकरण होता है।

  यदि कट लंबाई~$0$ है, तो कोई कट नहीं है और $\delta$ पहले से प्रसामान्य है।
  अतः मान लें कि~$\delta$ की कट कोटि और कट लंबाई दोनों~$> 0$ हैं। कोई सबसे
  ऊपर का, सबसे दायाँ महत्तम कट खंड चुनें। यदि कट खंड की लंबाई $>1$ है, तो
  क्रम-विनिमय रूपांतरण लगाएँ। यदि उसकी लंबाई~$1$ है, तो संगत अपचयन रूपांतरण
  लगाएँ। इससे ऐसा नया !!a{proof} मिलता है जिसकी कट कोटि समान पर कट लंबाई कम
  है, अथवा जिसकी कट कोटि कम है; इसलिए आगमन परिकल्पना लागू होती है।
\end{proof}

ऊपर के प्रमाण में प्रयुक्त युक्ति (सदा सबसे ऊपर के, सबसे दाएँ महत्तम कट से
निपटना) रूपांतरणों का ऐसा विशेष क्रम सुझाती है जिसमें उन्हें किसी !!{proof}
पर लगाने से प्रसामान्य !!{proof} प्राप्त होगा। आश्चर्यजनक रूप से क्रम का
वास्तव में कोई महत्त्व नहीं है। क्रम-विनिमय और अपचयन रूपांतरणों को किसी भी
क्रम में लगाने से अंततः प्रसामान्य प्रमाण मिलता है।

खंड और कट की परिभाषाएँ \Log{N2i} के लिए आसानी से सूत्रबद्ध की जा सकती हैं,
और \Log{N1i} के क्रम-विनिमय तथा अपचयन रूपांतरण सीधे \Log{N2i} में अनूदित
होते हैं। फलतः प्रसामान्यीकरण प्रमेय \Log{N2i} के लिए भी सत्य है।

\begin{thm}\ollabel{thm:normalization-N2i} प्रत्येक \Log{N2i}-!!{proof}
$\delta$ का प्रसामान्यीकरण होता है; अर्थात् उसे कट-विहीन
!!a{proof}~$\delta'$ में बदला जा सकता है।
\end{thm}

\end{document}
